% Manuscript skeleton for IEEE Transactions on Sustainable Energy (IEEEtran, journal mode).
%
% Numbers, tables and result figures are generated by `python -m windctrl.report` into
% tables/ and figures/. Anything not yet available is printed in red as [pending: ...];
% text the authors still have to write is marked with \todo{...}. Nothing in red is a result.
\documentclass[journal]{IEEEtran}

\usepackage{cite}
\usepackage{amsmath,amssymb}
\usepackage{graphicx}
\usepackage{xcolor}
\usepackage[hidelinks]{hyperref}

\newcommand{\pending}[1]{\textcolor{red}{[pending: #1]}}
\newcommand{\todo}[1]{\textcolor{red}{[to write: #1]}}
% Placeholder for a single value; \mbox makes it safe in both text and math mode.
\newcommand{\pv}[1]{\textcolor{red}{\mbox{[#1]}}}
\newcommand{\inputifexists}[2]{\IfFileExists{#1}{\input{#1}}{\par\noindent\pending{#2}\par}}
% \figifexists{file}{width}{placeholder text}
\newcommand{\figifexists}[3]{\IfFileExists{#1}{\includegraphics[width=#2]{#1}}{\fbox{\parbox{0.9\linewidth}{\centering\pending{#3}}}}}

% Measured values quoted in the text; red placeholders until the report has produced them.
\IfFileExists{tables/values.tex}{% Generated by windctrl.report; do not edit by hand.
\newcommand{\NormMs}{9047}
\newcommand{\NormVbar}{15.99}
\newcommand{\NormSigmaV}{2.91}
\newcommand{\NormSigmaW}{0.430}
\newcommand{\NormBetaBar}{0.234}
\newcommand{\NormSigmaBeta}{0.075}
\newcommand{\KanUpdates}{5}
\newcommand{\KanFailed}{0}
\newcommand{\MlpUpdates}{1}
\newcommand{\MlpFailed}{0}
}{}
\providecommand{\NormMs}{\pv{Ms}}
\providecommand{\NormVbar}{\pv{v-bar}}
\providecommand{\NormSigmaV}{\pv{sigma-v}}
\providecommand{\NormSigmaW}{\pv{sigma-omega}}
\providecommand{\NormBetaBar}{\pv{beta-bar}}
\providecommand{\NormSigmaBeta}{\pv{sigma-beta}}
\providecommand{\BestIpcGain}{\pv{gain}}
\providecommand{\KanUpdates}{\pv{updates}}
\providecommand{\KanFailed}{\pv{count}}
\providecommand{\LawOwn}{\pv{gs}}
\providecommand{\LawNext}{\pv{gn}}
\providecommand{\LawPrev}{\pv{gp}}
\providecommand{\LawSpeed}{\pv{k}}
\providecommand{\LawConst}{\pv{b}}
\providecommand{\GoneAbs}{\pv{G1 gain}}
\providecommand{\GonePhase}{\pv{G1 phase}}

\begin{document}

\title{Interpretable Individual Pitch Control of a 22-MW Offshore Wind Turbine With a
Kolmogorov--Arnold Reinforcement Learning Policy}

\author{\todo{author names and affiliations}%
\thanks{\todo{funding, manuscript dates, corresponding author}}}

\markboth{IEEE Transactions on Sustainable Energy}{}

\maketitle

\begin{abstract}
\todo{abstract, 150--250 words, once the results are final.} Individual pitch control (IPC) of
large offshore wind turbines has to reduce asymmetric blade loads without degrading speed and
power regulation. This paper develops a reinforcement-learning IPC law for the IEA 22-MW
reference turbine in which a Kolmogorov--Arnold network (KAN), readable by construction,
replaces the multilayer perceptron of a proximal policy optimization controller. The policy adds
one pitch offset per blade to the collective pitch command of the baseline controller.
\pending{headline results against PI and ROSCO IPC, and what the interpretation shows}
\end{abstract}

\begin{IEEEkeywords}
Fatigue, individual pitch control, interpretable machine learning, Kolmogorov--Arnold networks,
proximal policy optimization, reinforcement learning, wind turbines.
\end{IEEEkeywords}

% -------------------------------------------------------------------------------------------------
\section{Introduction}
\IEEEPARstart{O}{ffshore} wind turbines have reached rated powers above 20~MW, and the IEA
22-MW reference turbine has a rotor diameter of 284~m \cite{zahle2024iea22}. Above rated wind
speed, collective pitch control regulates rotor speed and power, whereas individual pitch control
(IPC) superimposes a pitch offset on each blade to attenuate the asymmetric blade loads, which are
dominated by the once-per-revolution (1P) component \cite{bossanyi2003ipc}. The Reference
Open-Source Controller (ROSCO) \cite{abbas2022rosco} includes an IPC module based on the Coleman
transformation and integral control.

Reinforcement learning (RL) has been applied to IPC \cite{coquelet2022rlipc}, but a policy
represented by a multilayer perceptron (MLP) does not reveal which input drives which blade.
Kolmogorov--Arnold networks (KANs) \cite{liu2025kan} replace scalar weights with learnable
one-dimensional functions, so a single KAN layer is an additive model whose terms can be inspected
directly.

\todo{literature review and positioning: recent pitch-control and RL work on large turbines,
post hoc explanation versus interpretable-by-construction policies, and how this paper differs
from earlier KAN-based IPC work on the 15-MW turbine.}

\todo{list of contributions, written after the results are final.}

% -------------------------------------------------------------------------------------------------
\section{Problem Formulation}
\subsection{Turbine, Baseline Controllers, and Simulation Interface}
The plant is the IEA 22-MW reference turbine on a monopile \cite{zahle2024iea22}, with a hub
height of 170~m and a rated rotor speed of 7.06~rpm (direct drive). It is simulated with OpenFAST
v4.2.1 \cite{jonkman2013fast}. The blades are modeled with the modal ElastoDyn module, and the
simulation time step is 0.002~s. ROSCO v2.10.5 \cite{abbas2022rosco} provides collective pitch and
torque control, and the pitch rate is limited to $\pm 2^\circ$/s, the value specified for this
turbine. Fig.~\ref{fig:framework} shows the overall framework.

The learned policy runs in a separate process and exchanges data with the simulation through the
ZeroMQ interface of ROSCO: every 0.02~s ROSCO sends its measurements and receives one pitch offset
per blade, which it adds to its own pitch command before the rate limit and the pitch actuator
model. The interface was extended to transmit the wind speed estimate and the collective pitch
angle.

Three controllers serve as baselines: PI, i.e., collective pitch control by ROSCO alone; ROSCO
IPC, with the Coleman transformation and integral control; and MLP-PPO, a policy with multilayer
perceptron networks trained in the same way as the proposed one. The integral gain of ROSCO IPC is
selected on the training winds (Section~\ref{sec:sweep}).

\begin{figure*}[!t]
\centering
\figifexists{figures/framework.pdf}{\textwidth}{framework figure}
\caption{Framework of the proposed interpretable KAN-PPO individual pitch controller. (a) Closed
loop: the KAN actor adds one pitch offset per blade to the collective pitch command of ROSCO
through the ZeroMQ interface. (b) Training with PPO on parallel OpenFAST rollouts. (c) Evaluation
on unseen winds and interpretation of the trained policy, whose explicit four-gain reduction is
executed in closed loop without retraining.}
\label{fig:framework}
\end{figure*}

\subsection{State, Action, and Reward}
Every 0.1~s, the policy receives the state
\begin{equation}
\mathbf{s} = \Big[\frac{\Delta M_1}{M_s},\ \frac{\Delta M_2}{M_s},\ \frac{\Delta M_3}{M_s},\
\frac{\hat v - \bar v}{\sigma_v},\ \frac{\omega_g - \omega_{\mathrm{ref}}}{\sigma_\omega},\
\frac{\beta - \bar\beta}{\sigma_\beta}\Big]
\label{eq:state}
\end{equation}
with $\Delta M_i = M_{y,i} - \tfrac{1}{3}\sum_j M_{y,j}$, where $M_{y,i}$ is the out-of-plane root
bending moment of blade $i$, $\hat v$ is the wind speed estimated by ROSCO, $\omega_g$ is the
generator speed, and $\beta$ is the collective pitch angle. The normalization constants are
measured from PI runs on the training winds: $M_s = \NormMs$~kN$\cdot$m (the standard deviation of
$\Delta M_i$), $\bar v = \NormVbar$~m/s, $\sigma_v = \NormSigmaV$~m/s,
$\omega_{\mathrm{ref}} = 7.06$~rpm, $\sigma_\omega = \NormSigmaW$~rpm,
$\bar\beta = \NormBetaBar$~rad, and $\sigma_\beta = \NormSigmaBeta$~rad.

The action $\mathbf{a} \in [-1,1]^3$ sets one offset per blade,
$\Delta\beta_i = \theta_{\max} a_i$ with $\theta_{\max} = 4^\circ$, which is smoothed by a
first-order filter with a time constant of 0.1~s. The reward is
\begin{equation}
\begin{aligned}
r = &-\frac{w_{\mathrm{bal}}}{3}\sum_i \frac{|\Delta M_i|}{M_s}
     - w_{\mathrm{pow}} \frac{|P - P_{\mathrm{rated}}|}{P_{\mathrm{rated}}} \\
    &-\frac{w_{\mathrm{smooth}}}{3}\sum_i |a_{i,k} - a_{i,k-1}|,
\end{aligned}
\label{eq:reward}
\end{equation}
where $P$ is the generator power, $P_{\mathrm{rated}} = 22$~MW, and $k$ is the decision index,
with $w_{\mathrm{bal}} = 1$, $w_{\mathrm{pow}} = 5$, and $w_{\mathrm{smooth}} = 0.1$. The policy
acts in the rotating frame: it uses neither an azimuth input nor a Coleman transformation.

\subsection{Simulation Protocol}
Wind fields are generated with TurbSim \cite{jonkman2009turbsim} using the IEC Kaimal model with
normal turbulence, class B \cite{iec61400}, at mean hub-height wind speeds of 13, 16, and 19~m/s.
Turbulence seeds 101 and 102 are used for training (300-s episodes) and seeds 201 and 202 for
evaluation (600-s runs). All runs include irregular waves with a significant wave height of
4.52~m and a peak period of 9.45~s. The first 60~s of each run are excluded from all statistics.
Damage-equivalent loads (DELs) are computed from the rainflow cycle counts \cite{astm1049} of
each load signal as
\begin{equation}
\mathrm{DEL} = \Big(\frac{1}{f_{\mathrm{eq}} T}\sum_k n_k S_k^m\Big)^{1/m},
\end{equation}
where $S_k$ and $n_k$ are the range and count of cycle $k$ (half counts for the residual), $T$ is
the analyzed duration, $f_{\mathrm{eq}} = 1$~Hz, and the W\"ohler exponent is $m = 10$ for the
blades and $m = 4$ for the tower. With 540~s per run, the DELs are short-term indicators rather
than lifetime estimates.

% -------------------------------------------------------------------------------------------------
\section{Kolmogorov--Arnold Policy and Its Interpretation}
\subsection{KAN Policy and Critic}
A single KAN layer maps $\mathbf{s}$ to the mean action,
\begin{equation}
\begin{aligned}
\mu_i &= \tanh z_i, \qquad z_i = b_i + \sum_{j=1}^{6} \varphi_{ij}(s_j), \\
\varphi_{ij}(x) &= w_{ij}x + \sum_{k=1}^{8} c_{ijk} B_k\big(\mathrm{clip}(x,-3,3)\big),
\end{aligned}
\label{eq:kan}
\end{equation}
where $B_k$ are the eight uniform cubic B-spline basis functions on five intervals over $[-3,3]$.
A separate layer of the same form, followed by a softplus, gives the standard deviation, so the
actor has 330 parameters. The critic is a KAN $[6,6,1]$ with 385 parameters whose output is scaled
by 10. The MLP-PPO baseline uses two hidden layers of 64 ReLU units for the actor (4998
parameters) and for the critic (4673 parameters).

Both policies are trained with PPO \cite{schulman2017ppo} and generalized advantage estimation
\cite{schulman2016gae}: learning rates of $3\times10^{-4}$ (actor) and $10^{-3}$ (critic) with
linear decay, $\gamma = 0.99$, $\lambda = 0.95$, clip range 0.2, 10 epochs per update, minibatch
size 512, no entropy bonus, gradient-norm clip 0.5, initial policy standard deviation 0.3, three
300-s episodes per update (one per mean wind speed, alternating the two training seeds), and 100
updates.

\subsection{Interpreting the Policy}
Before the $\tanh$, \eqref{eq:kan} is an exact sum, so the contribution of every input to every
blade is known at each instant. Since $\Delta M_1 + \Delta M_2 + \Delta M_3 = 0$, adding the same
slope to the three load edges of an output leaves the policy unchanged; this ambiguity is removed
by subtracting the mean load slope of each output and by centering every edge on the visited
states. The size of an edge is the standard deviation of $\theta_{\max}\varphi_{ij}$ over the
evaluation states.

Fitting a straight line to every edge turns the policy into a gain matrix, and averaging over the
blades gives a four-gain law,
\begin{equation}
\begin{aligned}
\Delta\beta_i = \theta_{\max}\tanh\Big(&\frac{g_s\Delta M_i + g_n\Delta M_{i+1} + g_p\Delta M_{i-1}}{M_s} \\
&+ k\,\frac{\omega_g - \omega_{\mathrm{ref}}}{\sigma_\omega} + b\Big),
\end{aligned}
\label{eq:law}
\end{equation}
which is executed in closed loop without retraining. The local gains combine into the 1P complex
gain $G_1 = g_s + g_n e^{+j120^\circ} + g_p e^{-j120^\circ}$: the offset of a blade equals its own
1P load scaled by $|G_1|$ and shifted by $\angle G_1$.

% -------------------------------------------------------------------------------------------------
\section{Results}
\subsection{Tuning of the ROSCO IPC Baseline}
\label{sec:sweep}
The integral gain of ROSCO IPC is swept on the training winds (seed 101), and the gain with the
lowest blade DEL is retained, $K_I = \BestIpcGain$~rad/(MN$\cdot$m$\cdot$s).
\inputifexists{tables/sweep.tex}{gain sweep table}
\todo{comment on the sweep, including any aborted runs and the behavior at the $2^\circ$/s rate
limit.}

\subsection{Training}
\begin{figure}[!t]
\centering
\figifexists{figures/training.pdf}{\columnwidth}{training curves}
\caption{Mean reward per decision during training (five-point moving mean), by mean wind speed.}
\label{fig:training}
\end{figure}
Fig.~\ref{fig:training} shows the training curves. The KAN policy completed \KanUpdates\ updates
with \KanFailed\ aborted episodes. \todo{compare the KAN and MLP curves and the reward level of
PI.}

\subsection{Closed-Loop Performance}
\inputifexists{tables/performance.tex}{performance table}
\inputifexists{tables/performance_by_wind.tex}{DELs per wind speed}
\begin{figure*}[!t]
\centering
\figifexists{figures/indicators.pdf}{\textwidth}{indicators per wind speed}
\caption{Indicators per mean wind speed and controller on the evaluation winds (bar: mean of the
two seeds; whisker: higher seed).}
\label{fig:indicators}
\end{figure*}
\todo{discussion of Table~\ref{tab:performance} and Fig.~\ref{fig:indicators}: KAN-PPO against PI,
ROSCO IPC, and MLP-PPO.}

% -------------------------------------------------------------------------------------------------
\section{Interpretability Analysis}
\begin{figure*}[!t]
\centering
\figifexists{figures/edges.pdf}{\textwidth}{KAN edge functions}
\caption{Edges of the trained policy in canonical form, $\theta_{\max}\varphi_{ij}$ versus the
physical input. Line styles distinguish the output blades. Shading marks the 1--99\% range of the
input on the evaluation winds.}
\label{fig:edges}
\end{figure*}
\inputifexists{tables/edge_sizes.tex}{edge sizes}
\inputifexists{tables/law_fit.tex}{fit of the simplified laws}
Near zero action, in degrees, MN$\cdot$m, and rpm, the four-gain law \eqref{eq:law} reads
\begin{equation}
\begin{aligned}
\Delta\beta_i \approx\ &\LawOwn\,\Delta M_i\ \LawNext\,\Delta M_{i+1}\ \LawPrev\,\Delta M_{i-1} \\
&\LawSpeed\,(\omega_g - \omega_{\mathrm{ref}})\ \LawConst,
\end{aligned}
\end{equation}
with $|G_1| = \GoneAbs^\circ$/(MN$\cdot$m) and $\angle G_1 = \GonePhase^\circ$.
\todo{interpretation of Fig.~\ref{fig:edges} and of the law; closed-loop comparison of the
four-gain law with the KAN (it appears in Table~\ref{tab:performance} when evaluated).}

% -------------------------------------------------------------------------------------------------
\section{Discussion}
\todo{discussion and limitations. Points that follow from the setup and should be stated: the
blades are modeled with ElastoDyn rather than BeamDyn; one training seed per policy and two
evaluation seeds per wind speed; fixed-bottom operation at 13--19~m/s in class B normal turbulence
and one sea state; short-term DELs.}

\section{Conclusion}
\todo{conclusion.}

\bibliographystyle{IEEEtran}
\bibliography{refs}

\end{document}
